%==============================================================================
% TCC — Rafael Alves Silva Rezende (V1 para orientador)
%==============================================================================
\documentclass[
    12pt,
    openright,
    oneside,
    a4paper,
    brazil
]{abntex2}

\usepackage[T1]{fontenc}
\usepackage[utf8]{inputenc}
\usepackage{amsmath}
\usepackage{amssymb}
\usepackage{graphicx}
\usepackage{indentfirst}
\usepackage{microtype}
\usepackage{booktabs}
\usepackage{tabularx}
\usepackage{hyperref}
\hypersetup{hidelinks}

% Macros de termos (evita conflito com glossaries/abntex2)
\newcommand{\termo}[1]{\csname termo@#1\endcsname}
\newcommand{\termo@llm}{LLM}
\newcommand{\termo@rag}{RAG}
\newcommand{\termo@ddl}{DDL}


\title{Arquitetura de Orquestração Cognitiva Baseada em LLMs: Integração de RAG e Text-to-SQL Aplicada a Dados de Transparência Municipal}
\author{Rafael Alves Silva Rezende}
\orientador{Prof. Dr. Denilson}
\instituicao{Universidade Federal de Lavras -- UFLA}
\local{Lavras -- MG}
\data{2026}
\tipotrabalho{Trabalho de Conclusão de Curso}

\begin{document}

\pretextual
\imprimircapa
\imprimirfolhaderosto

\begin{resumo}
Este trabalho apresenta o CIC (Centro de Informações Caetense), agente conversacional para consulta a dados públicos de Caeté (MG) via WhatsApp, integrando recuperação documental sobre o Diário Oficial e Text-to-SQL sobre bases SICOM/CGU com DuckDB. Descrevem-se a arquitetura de orquestração, técnicas de engenharia de contexto e avaliação formal com conjunto de referência de 80 consultas. Quatro experimentos compararam modelos e variantes de \textit{prompt}, revelando alta taxa de SQL válida e Execution Accuracy modesta sob comparação estrita de resultados.
\vspace{\onelineskip}
\textbf{Palavras-chave}: Inteligência Artificial; LLM; RAG; Text-to-SQL; Dados Abertos; Transparência Pública.
\end{resumo}

\tableofcontents
\textual

\chapter{Introdução}

\section{Contextualização}

Os dados públicos municipais constituem patrimônio informacional essencial à fiscalização social, à prestação de contas e à tomada de decisão democrática. No entanto, o acesso efetivo a essas informações permanece restrito a usuários com conhecimento técnico em portais de transparência, formatos tabulares e linguagens de consulta estruturada.

Nesse contexto, os Modelos de Linguagem de Grande Escala (\termo{llm}) permitem interfaces conversacionais capazes de interpretar perguntas em linguagem natural e traduzi-las em operações sobre bases governamentais. A combinação de Recuperação Aumentada por Geração (\termo{rag}) para documentos normativos e Text-to-SQL para consultas analíticas configura uma arquitetura híbrida com potencial de ampliar o acesso a dados de transparência.

Este trabalho apresenta o \textbf{CIC --- Centro de Informações Caetense}, agente conversacional voltado ao município de Caeté (MG), integrado ao WhatsApp por meio de fluxos automatizados. O sistema consulta dados públicos do SICOM e da CGU (2020--2025), além do Diário Oficial municipal via busca vetorial.

\section{Problema de pesquisa}

Apesar da abertura crescente de bases governamentais, usuários leigos enfrentam barreiras técnicas para formular consultas sobre despesas, receitas, folha de pagamento, licitações, contratos e programas sociais.

\section{Objetivo geral}

Desenvolver e avaliar empiricamente uma arquitetura de orquestração cognitiva baseada em \termo{llm}s, integrando \termo{rag} e Text-to-SQL para consultas conversacionais sobre dados de transparência municipal.

\section{Objetivos específicos}

\begin{itemize}
    \item Construir \textit{pipeline} de ingestão e normalização de dados SICOM/CGU para consulta analítica;
    \item Implementar subsistema Text-to-SQL com minificação de esquema e amostragem inteligente;
    \item Implementar subsistema de recuperação documental sobre o Diário Oficial;
    \item Orquestrar os fluxos em ambiente de produção via mensageria instantânea;
    \item Definir conjunto de referência e conduzir avaliação formal com métricas de execução, validade SQL, latência e custo;
    \item Realizar estudos de ablação comparando variantes de \textit{prompt} e modelos de linguagem.
\end{itemize}

\section{Justificativa}

A solução proposta alinha-se à governança digital e ao ODS~16, ao reduzir a barreira técnica de acesso à informação pública. A avaliação quantitativa em dados municipais reais contribui com evidências empíricas para decisões de arquitetura em soluções de tecnologia cívica (\textit{GovTech}).

\section{Organização do trabalho}

O Capítulo~2 apresenta o referencial teórico. O Capítulo~3 descreve metodologia e arquitetura. O Capítulo~4 analisa resultados. O Capítulo~5 discute limitações e implicações. O Capítulo~6 encerra com conclusões.

\chapter{Referencial Teórico}

\section{Modelos de linguagem e dados governamentais}

Os Modelos de Linguagem de Grande Escala (\termo{llm}) são capazes de generalizar padrões linguísticos e realizar raciocínio aproximado sobre instruções em linguagem natural. Em bases governamentais heterogêneas, o desafio central é condicionar o modelo com contexto suficiente sem exceder limites de tokens e custo operacional.

\section{Recuperação aumentada por geração}

A técnica \termo{rag} complementa o \termo{llm} com recuperação de trechos documentais relevantes antes da geração da resposta, sendo adequada a perguntas sobre leis, decretos e normas municipais.

\section{Text-to-SQL}

Text-to-SQL converte perguntas em consultas estruturadas sobre bancos relacionais, permitindo respostas quantitativas sobre despesas, receitas, folha e licitações. A literatura emprega métricas como \textit{Execution Accuracy} e \textit{Valid SQL Rate} para avaliação.

\section{Orquestração cognitiva}

Designa a camada que interpreta a intenção do usuário e seleciona dinamicamente o mecanismo mais adequado. Plataformas de automação permitem integrar APIs, modelos e canais conversacionais em arquitetura modular.

\section{Trabalhos relacionados}

Benchmarks Text-to-SQL (Spider, BIRD) estabelecem protocolos com conjuntos de referência e comparação de resultados executados. Este trabalho adapta esses princípios ao contexto municipal brasileiro.

\chapter{Metodologia}

\section{Tipo de pesquisa}

Trata-se de pesquisa aplicada, de natureza experimental e desenvolvimento tecnológico, com validação funcional em ambiente real e avaliação quantitativa por \textit{benchmark} reprodutível em ambiente controlado.

\section{Contexto e dados}

O município-alvo é Caeté (MG), código IBGE/SICOM 3110004. As fontes incluem:

\begin{itemize}
    \item \textbf{SICOM}: receitas, despesas, empenhos, folha, frota, licitações e contratos (2020--2025);
    \item \textbf{CGU}: transferências federais e programas sociais cruzáveis por documento do credor;
    \item \textbf{Diário Oficial}: documentos normativos indexados em Qdrant para recuperação aumentada.
\end{itemize}

Um processo ETL em Python consolida os arquivos tabulares, gera metadados de esquema (\termo{ddl}) e expõe uma API analítica com DuckDB \textit{in-memory}.

\section{Arquitetura do sistema}

A Figura~\ref{fig:arquitetura} (Apêndice) resume o fluxo conversacional:

\begin{enumerate}
    \item O usuário envia a pergunta via WhatsApp;
    \item A plataforma de automação recebe o evento e aciona o agente orquestrador (GPT-4o);
    \item O agente seleciona recuperação documental (Qdrant) ou tradução Text-to-SQL;
    \item O subsistema SQL utiliza GPT-4.1-mini com \textit{prompt} contendo esquema compactado e amostras de dados;
    \item O resultado tabular é interpretado pelo orquestrador e devolvido em linguagem natural ao usuário.
\end{enumerate}

\section{Engenharia de contexto}

\subsection{Minificação de esquema (\textit{DDL minification})}

Consiste na remoção de redundâncias, padronização de descrições e simplificação do esquema relacional enviado ao modelo, reduzindo tokens sem eliminar restrições semânticas críticas.

\subsection{Amostragem inteligente (\textit{smart sampling})}

Inclusão de exemplos reais de valores por coluna relevante, visando melhorar inferência de filtros, \textit{joins} e convenções locais.

\subsection{Variantes de \textit{prompt} avaliadas}

Para o estudo de ablação, foram definidas duas variantes de contexto:

\begin{itemize}
    \item \textbf{Prompt compacto} ($\approx$16\,500 tokens): esquema minificado e amostras --- configuração adotada em produção;
    \item \textbf{Prompt estendido} ($\approx$35\,000 tokens): esquema com descrições ampliadas --- hipótese de maior cobertura semântica.
\end{itemize}

\section{Conjunto de referência (\textit{golden dataset})}

Foi construído um conjunto de 80 consultas em linguagem natural, estratificadas por dificuldade (fácil, médio, difícil) e por domínio temático (recursos humanos, frota, licitações, cruzamentos analíticos e regras especiais). Cada item contém a pergunta do usuário, a consulta SQL de referência validada no DuckDB e metadados para métricas auxiliares.

Dez consultas-base foram expandidas com variações paramétricas (anos, limites e agregações), totalizando 80 casos com SQL executável.

\section{Protocolo de avaliação}

A avaliação Text-to-SQL foi conduzida por um pipeline automatizado, desacoplado do canal WhatsApp, garantindo reprodutibilidade e controle de orçamento de API. O protocolo registra, por consulta: SQL gerada, resultados executados, métricas, latência, consumo de tokens e custo estimado.

\subsection{Métricas de desempenho}

As métricas adotadas seguem a literatura de Text-to-SQL e foram complementadas por indicadores operacionais:

\begin{description}
    \item[EX --- \textit{Execution Accuracy}] Proporção de consultas em que o resultado da SQL gerada coincide com o da referência (comparação de multiconjunto de linhas normalizadas). Mede se o cidadão receberia os mesmos valores numéricos da consulta validada.
    \item[VSR --- \textit{Valid SQL Rate}] Proporção de SQLs executadas sem erro no DuckDB. Distingue falha sintática de falha semântica.
    \item[NEA --- \textit{Non-Empty Answer}] Indica se a consulta retornou dados quando o conjunto de referência espera resultado não vazio.
    \item[TSA --- \textit{Table Selection Accuracy}] Verifica se as tabelas esperadas foram utilizadas na SQL gerada.
    \item[CHS --- \textit{Condition Heuristic Score}] Mede a presença de filtros relevantes (ano, órgão, categoria) na consulta produzida.
\end{description}

Em termos interpretativos, o padrão \textit{VSR elevado com EX moderado} indica que o modelo domina a sintaxe e o esquema relacional, mas frequentemente interpreta a pergunta de forma distinta da referência manual --- limitação conhecida da métrica EX estrita, que não admite consultas semanticamente equivalentes com formulação diferente.

\subsection{Métricas operacionais}

Registraram-se latência de geração, latência de execução SQL, tokens de entrada e saída e custo estimado em dólares, conforme tabela de preços vigente da API OpenAI.

\subsection{Experimentos conduzidos}

\begin{enumerate}
    \item \textbf{E1 --- Baseline}: GPT-4.1-mini com \textit{prompt} compacto;
    \item \textbf{E2 --- Ablação de contexto}: GPT-4.1-mini com \textit{prompt} estendido;
    \item \textbf{E3 --- Ablação de modelo (econômico)}: GPT-4o-mini com \textit{prompt} compacto;
    \item \textbf{E4 --- Ablação de modelo (avançado)}: GPT-4o com \textit{prompt} compacto.
\end{enumerate}

O orçamento total de experimentação foi limitado a aproximadamente US\$~14,00, com margem de segurança para interrupção automática.

\section{Ambiente experimental}

\begin{itemize}
    \item API analítica DuckDB em ambiente local;
    \item Modelos OpenAI via API oficial, temperatura 0,0;
    \item Saída estruturada em JSON contendo a consulta SQL e justificativa técnica;
    \item Persistência incremental de resultados para retomada em caso de interrupção.
\end{itemize}

\chapter{Análise de Resultados}

\section{Configuração baseline (E1)}

O experimento E1 avaliou 80 consultas com GPT-4.1-mini e \textit{prompt} compacto ($\approx$16\,500 tokens de contexto). A Tabela~\ref{tab:baseline} resume as métricas agregadas.

\begin{table}[htbp]
\centering
\caption{Resultados do experimento E1 --- baseline ($n=80$)}
\label{tab:baseline}
\begin{tabular}{lr}
\toprule
\textbf{Métrica} & \textbf{Valor} \\
\midrule
Execution Accuracy (EX) & 11,25\% (9/80) \\
Valid SQL Rate (VSR) & 97,50\% (78/80) \\
Custo total estimado & US\$ 0,56 \\
Latência média de geração & 6,2 s \\
Latência média de execução SQL & 1,1 s \\
\bottomrule
\end{tabular}
\end{table}

\section{Resultados por dificuldade}

\begin{table}[htbp]
\centering
\caption{EX e VSR por nível de dificuldade (E1)}
\label{tab:diff}
\begin{tabular}{lrrr}
\toprule
\textbf{Dificuldade} & \textbf{$n$} & \textbf{EX (\%)} & \textbf{VSR (\%)} \\
\midrule
Fácil & 22 & 13,64 & 100,00 \\
Médio & 44 & 9,09 & 97,73 \\
Difícil & 14 & 14,29 & 92,86 \\
\bottomrule
\end{tabular}
\end{table}

\section{Resultados por categoria temática}

As categorias com maior EX relativo foram \textit{Cálculo Fiscal} (30,77\%) e \textit{Consulta Simples} (16,67\%). Domínios como \textit{RH/Folha} e \textit{Despesas e Pagamentos} apresentaram EX nulo no baseline, embora o VSR permanecesse elevado --- indicando consultas sintaticamente válidas, porém semanticamente divergentes da referência.

\section{Interpretação}

O alto VSR (97,5\%) combinado ao EX modesto (11,25\%) sugere que a principal limitação observada não está na geração de SQL executável, mas na equivalência estrita entre a intenção da pergunta e a consulta produzida. Em diversos casos, o modelo retornou o mesmo número de linhas da referência com agregações ou filtros distintos.

\section{Comparativo de ablações}

A Tabela~\ref{tab:ablation} consolida os quatro experimentos. O GPT-4o (E4) obteve maior EX (18,75\%) e VSR equivalente ao baseline (97,5\%), porém com custo aproximadamente seis vezes superior (US\$~3,45) e latência média de geração de 29,5~s, o que limita sua adoção direta em canal conversacional síncrono.

\begin{table}[htbp]
\centering
\caption{Comparativo dos experimentos de ablação ($n=80$)}
\label{tab:ablation}
\small
\begin{tabular}{p{4.8cm}rrrr}
\toprule
\textbf{Experimento} & \textbf{EX (\%)} & \textbf{VSR (\%)} & \textbf{Custo (US\$)} & \textbf{Lat. (s)} \\
\midrule
E1 --- Baseline (GPT-4.1-mini, prompt compacto) & 11,25 & 97,50 & 0,56 & 6,2 \\
E2 --- Contexto estendido (GPT-4.1-mini) & 15,00 & 96,25 & 1,16 & 8,2 \\
E3 --- Modelo econômico (GPT-4o-mini) & 16,25 & 90,00 & 0,21 & 4,0 \\
E4 --- Modelo avançado (GPT-4o) & 18,75 & 97,50 & 3,45 & 29,5 \\
\bottomrule
\end{tabular}
\end{table}

\section{Validação funcional}

Além do \textit{benchmark} offline, o sistema foi validado em ambiente operacional via WhatsApp, com consultas reais de usuários, confirmando a viabilidade da orquestração integrada entre automação, mensageria e subsistemas de recuperação documental e consulta analítica.
\section{Avaliação da resposta em linguagem natural}

Além das métricas sobre a consulta SQL, foi conduzida avaliação da \textbf{cadeia completa} do subsistema analítico: execução da consulta, geração da resposta ao cidadão pelo agente orquestrador e julgamento automático por outro modelo de linguagem (LLM-as-Judge).

\subsection{Protocolo}

Para cada item amostrado do experimento E1, o protocolo registrou: (i)~pergunta do usuário e SQL gerada; (ii)~resultado tabular executado no DuckDB; (iii)~resposta em linguagem natural produzida pelo orquestrador (GPT-4o); (iv)~avaliação do juiz (GPT-4.1-mini) em quatro dimensões Likert (adequação semântica, fidelidade aos dados, completude e clareza).

A métrica \textbf{JAR} (\textit{Judge Agreement Rate}) corresponde à proporção de respostas classificadas como \textit{corretas} ou \textit{parcialmente corretas} pelo juiz.

\subsection{Piloto ($n=5$)}

O piloto inicial confirmou a viabilidade do pipeline. Observou-se JAR de 100\% com EX de 0\%, reforçando que a métrica EX estrita subestima a utilidade percebida quando a resposta natural atende à pergunta com os dados retornados.

\subsection{Resultados estratificados ($n=30$)}

A amostra estratificada compreendeu 10 consultas fáceis, 10 médias e 10 difíceis. A Tabela~\ref{tab:jar} resume os resultados.

\begin{table}[htbp]
\centering
\caption{Avaliação LLM-as-Judge sobre respostas naturais (E1, $n=30$)}
\label{tab:jar}
\begin{tabular}{lrr}
\toprule
\textbf{Métrica} & \textbf{Valor} & \textbf{Observação} \\
\midrule
JAR & 100,00\% (30/30) & Respostas corretas ou parcialmente corretas \\
EX (mesma amostra) & 20,00\% (6/30) & Equivalência estrita SQL \\
Adequação semântica (média) & 4,87 / 5 & Escala Likert do juiz \\
Fidelidade aos dados (média) & 4,87 / 5 & \\
Completude (média) & 4,73 / 5 & \\
Clareza ao cidadão (média) & 4,77 / 5 & \\
Custo total da amostra & US\$ 0,30 & Orquestrador + juiz \\
\bottomrule
\end{tabular}
\end{table}

\begin{table}[htbp]
\centering
\caption{JAR e EX por dificuldade (amostra $n=30$)}
\label{tab:jar-diff}
\begin{tabular}{lrrr}
\toprule
\textbf{Dificuldade} & \textbf{$n$} & \textbf{JAR (\%)} & \textbf{EX (\%)} \\
\midrule
Fácil & 10 & 100,00 & 20,00 \\
Médio & 10 & 100,00 & 20,00 \\
Difícil & 10 & 100,00 & 20,00 \\
\bottomrule
\end{tabular}
\end{table}

\subsection{Interpretação}

O contraste entre EX (20\%) e JAR (100\%) na mesma amostra indica que, embora a formulação SQL frequentemente difira da referência manual, o orquestrador consegue comunicar resultados úteis e fidedignos aos dados executados. Essa evidência complementa as limitações da métrica EX e justifica a avaliação em duas camadas: sintaxe/resultado tabular e adequação da resposta ao usuário final.

Os registros completos de entrada e saída de cada etapa foram arquivados para auditoria e reprodutibilidade do experimento.

\chapter{Discussão}

\section{Validade sintática versus acurácia de execução}

Os quatro experimentos convergem para um padrão consistente: VSR entre 90\% e 97\% e EX entre 11\% e 19\%. O modelo demonstra domínio da sintaxe SQL e do mapeamento esquema-coluna, mas produz com frequência consultas semanticamente plausíveis que divergem da referência anotada manualmente.

Esse comportamento é compatível com a literatura de Text-to-SQL em bases reais, nas quais múltiplas formulações SQL podem responder adequadamente à mesma pergunta, embora apenas uma tenha sido registrada como gabarito.

\section{Impacto do tamanho do contexto}

O \textit{prompt} estendido (E2) apresentou EX ligeiramente superior ao compacto (15,0\% \textit{versus} 11,25\%), porém com custo aproximadamente duplicado e latência maior. Para o cenário municipal com orçamento limitado, a variante compacta permanece mais custo-efetiva, dado o ganho marginal observado.

\section{Comparação entre modelos}

O GPT-4o-mini (E3) alcançou EX de 16,25\% com o menor custo (US\$~0,21), porém com queda de VSR para 90\%. O GPT-4o (E4) obteve o melhor EX (18,75\%), com destaque em consultas difíceis (42,86\%), à custa de latência e preço incompatíveis com respostas interativas imediatas.

\section{Limitações}

\begin{itemize}
    \item Lacuna de dados de folha de pagamento e cadastro de servidores para 2020--2022, impactando consultas de recursos humanos;
    \item Referências do conjunto \textit{golden} podem não abranger todas as SQLs semanticamente corretas;
    \item A avaliação offline não mede diretamente a qualidade da resposta em linguagem natural ao usuário final;
    \item As métricas reportadas concentram-se no subsistema Text-to-SQL; a avaliação documental (RAG) permanece como extensão futura.
\end{itemize}

\section{Implicações práticas}

Apesar do EX estrito modesto, o sistema em produção demonstra utilidade para cidadãos e jornalistas ao combinar consultas válidas, explicação técnica e interface conversacional. Os resultados subsidiam decisões de engenharia quanto a modelo, tamanho de contexto e orçamento operacional.

\chapter{Conclusão}

Este trabalho apresentou o CIC, arquitetura híbrida de orquestração cognitiva para consultas sobre dados de transparência municipal de Caeté (MG), integrando recuperação documental e Text-to-SQL via DuckDB.

A avaliação com conjunto de referência de 80 consultas demonstrou que o subsistema Text-to-SQL com GPT-4.1-mini e esquema minificado alcança elevada taxa de SQL válida (VSR), porém Execution Accuracy modesta sob métrica estrita de equivalência de resultados.

As contribuições incluem: (i)~arquitetura replicável para GovTech municipal; (ii)~\textit{pipeline} ETL e API analítica; (iii)~protocolo de avaliação reprodutível; e (iv)~evidências empíricas para escolha de modelo e tamanho de contexto.

Como trabalhos futuros, propõe-se: revisão humana das referências do conjunto \textit{golden}; avaliação da qualidade das respostas em linguagem natural; testes com usuários reais; e autocorreção iterativa de consultas SQL.


\appendix
\chapter{Arquitetura do Sistema}
\label{ap:arquitetura}

\begin{figure}[htbp]
\centering
\fbox{\parbox{0.92\textwidth}{
\small
\textbf{Fluxo do CIC:}\quad
WhatsApp $\rightarrow$ Gateway de mensageria $\rightarrow$ Plataforma de automação $\rightarrow$ Agente orquestrador\\
$\rightarrow$ \{Recuperação documental (vetorial) \textbf{ou} Text-to-SQL (modelo dedicado $\rightarrow$ API analítica)\}\\
$\rightarrow$ Resposta em linguagem natural $\rightarrow$ WhatsApp
}}
\caption{Visão geral da arquitetura de orquestração cognitiva}
\label{fig:arquitetura}
\end{figure}

\end{document}
